\section*{Resumen ejecutivo}
La productividad puede aumentar porque las empresas mejoran o porque cambia la distribución del empleo entre ellas. Distinguir ambos mecanismos permite entender si el avance observado es generalizado o si depende de que ganen peso las empresas más productivas y salgan de la muestra las menos productivas.

Este boletín aplica la descomposición dinámica de Olley--Pakes de \citet{melitz2015} a las encuestas anuales de industria (EAM), comercio (EAC) y servicios (EAS) del DANE. La variable empresarial es el logaritmo del valor agregado real por persona ocupada, ponderado por la participación de cada empresa en el empleo. Los resultados corresponden a las empresas cubiertas por estas encuestas; no representan toda la economía colombiana.

\begin{mdframed}[backgroundcolor=blue!4,linecolor=blue!25]
\textbf{Tres hallazgos para la discusión}
\begin{enumerate}[leftmargin=*,itemsep=5pt]
\item \textbf{Industria: el avance se concentra en la composición empresarial.} En 2019--2024, el índice agregado aumenta 10,4 puntos logarítmicos. La reasignación entre sobrevivientes aporta 7,5 y la salida 3,8; la media sin ponderar de las sobrevivientes aporta $-0,1$ y la entrada $-0,7$.
\item \textbf{Servicios: también mejoran las empresas que permanecen.} En 2019--2024, el aumento es de 13,4 puntos: 8,4 provienen del cambio de la media de sobrevivientes y 6,9 de la covarianza entre tamaño y productividad. La entrada resta 3,7.
\item \textbf{Comercio: la lectura depende del periodo y de la cobertura.} El índice cae 0,9 puntos en 2015--2019 y aumenta 8,7 en 2020--2024. La ventana reciente comienza en pandemia porque el identificador empresarial se rompe entre 2019 y 2020.
\end{enumerate}
\end{mdframed}

Estos hallazgos justifican estudiar simultáneamente la mejora de las empresas existentes y las condiciones para que las más productivas crezcan. No permiten concluir que los cierres sean deseables, que la entrada perjudique la economía o que una política específica haya causado los cambios.

\textbf{Alcance de esta versión.} Se presentan resultados descriptivos iniciales, con identidades contables verificadas. Quedan pendientes pruebas con deflactores por actividad, valores extremos, umbrales de tamaño y definiciones más estrictas de entrada y salida. Las cifras son cambios acumulados entre extremos, no tasas anuales. Este boletín complementa el informe agregado del DANE disponible en \texttt{main.tex}, sin sustituirlo.
\clearpage
